\section*{Introduction}

Modern deep learning architectures (e.g BERT, Vision Transformers) tend to have billions of parameters and are trained on expensive
and resource-intensive GPU clusters. Thus, the capacity need of storing these models and the latency their inference induces become bottlenecks of
deploying them in resource-constrained environments and latency-sensitive applications.  Besides designing more efficient architectures and training 
them from scratch, there are several well-established techniques of compressing a given model after training, such as network pruning, quantization,
knowledge distillation and tensor decomposition. 

My work focuses on network pruning which involves removing weights or groups of weights from the 
network whose contribution to the accuracy of the model is considered the smallest. I examined the effect of using three different regularization 
terms (L0, L1, weighted L1) in the objective function, and compared their capability of training a sparse network which retains accuracy when pruned 
to various degrees. 